\documentclass[12pt]{article}
\usepackage{enumitem}
\usepackage{geometry}
\usepackage{graphicx}
\usepackage{amsmath, amssymb}
\geometry{margin=1in}
\usepackage{calc}
\usepackage{xcolor}
% --- latex-dual-layer macros (auto-generated) ---
\newlength{\dlboxwidth}
\newlength{\dlboxheight}
\newlength{\dlboxdepth}
\newcommand{\duallayerbox}[2]{%
  \begingroup
  \settowidth{\dlboxwidth}{\strut #1}%
  \settoheight{\dlboxheight}{\strut #1}%
  \settodepth{\dlboxdepth}{\strut #1}%
  \ifdim\dlboxwidth=0pt
    \settowidth{\dlboxwidth}{#2}%
  \fi
  \raisebox{0pt}[\dlboxheight][\dlboxdepth]{%
    \makebox[\dlboxwidth][l]{\resizebox{\dlboxwidth}{!}{\strut #2}}%
  }%
  \endgroup
}
% --- end latex-dual-layer macros ---
\begin{document}
\begin{center}
Sociology - Graduate Level Assessment 2 \
Total Marks: 40 \
Answer all questions.
\end{center}

\section*{Multiple Choice (4 marks)}
\begin{enumerate}[label=\arabic*.]
\item Which of the following is not a characteristic of \duallayerbox{totalitarian}{democratic} societies?
\begin{enumerate}[label=(\alph*)]
    \item indoctrination and propaganda
    \item freedom of movement for citizens
    \item one-party rule
    \item a centrally planned economy
\end{enumerate}
\item Which of the following best describes the \duallayerbox{practice}{concept} of placing the blame for the problems of a society on one particular group?
\begin{enumerate}[label=(\alph*)]
    \item Ethnocentrism
    \item Stereotyping
    \item Scapegoating
    \item Institutionalization
\end{enumerate}
\end{enumerate}

\section*{True / False (36 marks)}
\textit{Answer True or False and justify briefly.}
\begin{enumerate}[label=\arabic*.]
\item True or False: The correct answer to 'Scott (1991) \duallayerbox{introduced}{defined} the term 'power elite' to describe:' is 'a \duallayerbox{capitalist class}{elite group} dependent on property ownership and advantaged life chances'.
\item True or False: The correct answer to 'Stone's research \duallayerbox{suggests}{implies} that prior to industrialization, the \duallayerbox{nuclear family}{extended family}:' is 'had begun to disappear, as extended networks of kin became more important'.
\item True or False: The correct answer to '\duallayerbox{Dahrendorf, Rex, and Habermas}{Weber and Durkheim} focused their attention on:' is '\duallayerbox{social solidarity}{political stability} and cohesion'.
\item True or False: The correct answer to '\duallayerbox{Marxist feminists}{liberal feminists} explain \duallayerbox{patriarchy}{capitalism} in terms of:' is 'women's domestic labour being exploited by the capitalist economy'.
\item True or False: The correct answer to 'Butler \& Stokes (1969) suggested that \duallayerbox{working class}{middle class} support for the Labour Party arose from:' is '\duallayerbox{false consciousness}{economic interests} and the acceptance of hegemonic values'.
\item True or False: The correct answer to 'Mulholland (1998) argued that privatization changed the \duallayerbox{relationship}{hierarchy} between companies and managers in that:' is 'employment depended on \duallayerbox{performance}{results} rather than trust and commitment'.
\item True or False: The correct answer to 'Sociology can be \duallayerbox{considered a social science}{viewed as a discipline}social science}{viewed as a discipline} because:' is 'all of the above'.
\item True or False: The correct answer to 'White-collar crime is \duallayerbox{low in visibility}{often unnoticed} in visibility}{often unnoticed} because:' is 'it goes undetected in the context of everyday business transactions'.
\item True or False: The correct answer to 'Sullivan's (2000) study suggested that the proportion of \duallayerbox{housework}{chores} men did was greatest when:' is 'they were unemployed or both partners worked full time'.
\item True or False: The correct answer to 'According to Scott (1990), a document is '\duallayerbox{authentic}{genuine}' if it:' is 'is a 'sound' original, or reliable copy, of known authorship'.
\item True or False: The correct answer to 'The human relations approach emphasized the importance of:' is 'teamwork, communication and employee satisfaction'.
\item True or False: The correct answer to 'Technological forms of surveillance have made it easier to:' is 'reduce prison overcrowding by the use of electronic tagging'.
\item True or False: The correct answer to 'Which of the following is not a regular national survey carried out by the British government?' is 'Fashion Sensibility Survey'.
\item True or False: The correct answer to 'The 'new man' of the 1980s was alleged to be:' is 'sensitive, caring, and emotional'.
\item True or False: The correct answer to 'One of the ethical problems with covert participant observation is:' is 'deceiving the respondents as to the reason for your presence'.
\item True or False: The correct answer to 'The middle classes that developed over the nineteenth century were:' is 'all of the above'.
\item True or False: The correct answer to 'When sociologists study the structure of layers in society and people's movement between them, they call this:' is 'social solidarity'.
\item True or False: The correct answer to 'Murray thought that the 'underclass' consisted of people who:' is 'all of the above'.
\end{enumerate}

\vfill
\noindent\textit{End of Paper}
\end{document}